%% kapa: exam=1 round=25 year=2014 period=2 form=A subject=appraisal_law subject_ko=부동산관계법규 seq=26 number=26
%% source: Q-Net 감정평가사 기출문제 (cst00309) artlSeq=5101207 file=extracted/1008/5101207/2014년 제25회 감정평가사 1차시험 기출문제/2교시(A형).hwp
%% conversion: kordoc 4.24.51 -> tools/gen1.py (auto); review=proofread-sonnet; flags=none
\begin{question}{26}
국유재산법령상 국유재산관리기금에 관한 설명으로 옳지 않은 것은?
\begin{choices}
  \item 국유재산관리기금은 국유재산의 취득에 필요한 비용의 지출에 사용할 수 있다.
  \item 국유재산관리기금은 총괄청 소관 일반재산의 관리ㆍ처분에 필요한 비용의 지출에 사용할 수 있다.
  \item 금융회사 등으로부터의 차입금은 국유재산관리기금의 재원에 해당한다.
  \item 총괄청은 국유재산관리기금의 결산보고서 작성에 관한 사무를 한국자산관리공사에 위탁한다.
  \item 국유재산관리기금에서 취득한 재산은 특별회계 소속으로 한다.
\end{choices}
\end{question}
\answer{5}
